% ****** Start of file apssamp.tex ******
%
%   This file is part of the APS files in the REVTeX 4.2 distribution.
%   Version 4.2a of REVTeX, December 2014
%
%   Copyright (c) 2014 The American Physical Society.
%
%   See the REVTeX 4 README file for restrictions and more information.
%
% TeX'ing this file requires that you have AMS-LaTeX 2.0 installed
% as well as the rest of the prerequisites for REVTeX 4.2
%
% See the REVTeX 4 README file
% It also requires running BibTeX. The commands are as follows:
%
%  1)  latex apssamp.tex
%  2)  bibtex apssamp
%  3)  latex apssamp.tex
%  4)  latex apssamp.tex
%
\documentclass[%
 reprint,
%superscriptaddress,
%groupedaddress,
%unsortedaddress,
%runinaddress,
%frontmatterverbose, 
%preprint,
%preprintnumbers,
%nofootinbib,
%nobibnotes,
%bibnotes,
 amsmath,amssymb,
 aps,
pra,
% prb,
%rmp,
%prstab,
%prstper,
%floatfix,
]{revtex4-2}

\usepackage{braket}
\usepackage{hyperref}
\hypersetup{
    colorlinks=true,
    linkcolor=black,
    citecolor=black,
    urlcolor=black,
    }
\usepackage{graphicx}% Include figure files
\usepackage{dcolumn}% Align table columns on decimal point
\usepackage{bm}% bold math
%\usepackage{hyperref}% add hypertext capabilities
%\usepackage[mathlines]{lineno}% Enable numbering of text and display math
%\linenumbers\relax % Commence numbering lines

%\usepackage[showframe,%Uncomment any one of the following lines to test 
%%scale=0.7, marginratio={1:1, 2:3}, ignoreall,% default settings
%%text={7in,10in},centering,
%%margin=1.5in,
%%total={6.5in,8.75in}, top=1.2in, left=0.9in, includefoot,
%%height=10in,a5paper,hmargin={3cm,0.8in},
%]{geometry}

\begin{document}

\title{Self-Interaction-Free Levy--Perdew--Sahni Density Functional Theory}% Force line breaks with \\

\author{Ivan P. Bosko and Viktor N. Staroverov}
 % \email{vstarove@uwo.ca}
\affiliation{%
 Department of Chemistry, The University of Western Ontario, London, Ontario N6A 5B7, Canada
}%

\date{\today}

\begin{abstract}

\end{abstract}
\maketitle

\section{Introduction}

Levi, Perdew, and Sahni derived a Schrödinger-like equation for a square-root of the electron density
\begin{equation}\label{eq: LPS Differential Equation}
    \left\{
    -\frac{1}{2}\nabla^2
    +
    v_{\text{eff}}([n];\textbf{r})
    \right\}
    n^{1/2}(\textbf{r})
    =
    \mu n^{1/2}(\textbf{r}),
\end{equation} 
where $v_{\text{eff}}([n];\textbf{r})$ acts as the multiplicative effective potential. The eigenvalue $\mu$ represents the chemical potential, which corresponds to the negative of the ionization energy in exact DFT~\cite{Levy.1984.PRA.30.2745}. The exact form of the effective potential has to be approximated 
\begin{equation}\label{eq: LPS Effective Potential}
    v_{\text{eff}}([n]; \textbf{r})
    =
    v_{\text{ext}}(\textbf{r})
    +
    \int\frac{n(\textbf{r}')}{|\textbf{r}-\textbf{r}'|}
    d\textbf{r}'
    +
    \frac{\delta G[n]}{\delta n(\textbf{r})}.
\end{equation}
Acharya \textit{et al.}~\cite{Acharya.1980.PNAS.77.6978} as well as Cedillo \textit{et al.}~\cite{Cedillo.1986.JCP.85.7188} explored the possibility for including the $N$-dependent factor into the approximate functionals to improve their performance. Here, we elaborate on this idea through analytical derivation of $N$-dependent self-interaction-corrected $G[n]$ functionals in LPS theory. Our results suggest a universal way of developing such approximations without empirical fitting and with straightforward generalization to spin-polarized systems.

In one-determinantal approximations (Hartree--Fock, Kohn--Sham DFT), the exact exchange, $K[\{\varphi_i\}]$, is given by
\begin{equation}
    K[\{\varphi_i\}]
    =
    \sum^N_{i,j}
    K_{ij}
    =
    \sum^N_{i=1}
    K_{ii}
    +
    \sum^N_{i\neq j}
    K_{ij},
\end{equation}
where the sum over $i$ only represents the self-exchange and the sum over all distinct pairs of $i$ and $j$ represents the interelectronic exchange contribution.
Each exchange integral, $K_{ij}$, is expressed via spin-orbitals as
\begin{equation}
    K_{ij}
    =
    -
    \int\int
    \frac{\varphi_i^*(\textbf{x})\varphi_j(\textbf{x})
          \varphi_j^*(\textbf{x}')\varphi_i(\textbf{x}')}
          {|\textbf{x}-\textbf{x}'|}
    d\textbf{x}d\textbf{x}'.
\end{equation} 
For one-orbital systems, the whole exact exchange reduces to the self-exchange alone and is equal to the Fermi--Amaldi functional
\begin{equation}
    K_{\text{FA}}[n]
    =
    -
    \frac{1}{2N}
    \int\int
    \frac{n(\textbf{r})n(\textbf{r}')}
    {|\textbf{r}-\textbf{r}'|}
    d\textbf{r}d\textbf{r}'.
\end{equation}
Recently, we have shown that the Fermi--Amaldi functional is the exact exchange in LPS DFT. Due to the nature of the boson-like LPS reference system, the role of the LPS exact exchange (FA) is reduced to the cancelation of self-repulsion between coulombically-interacting particles, i.e., self-exchange~\cite{Bosko.2023.PRA.108.052803}. When used as a full electronic exchange within LPS theory, FA can deliver better performance than LDA and PBE, given it is combined with a proper kinetic energy functional. However, to better describe interacting electrons using the boson-like reference system, it seems reasonable to approximate the interelectronic exchange by an approximate exchange functional.

\section{Self-Interaction-Corrected LPS Theory}

\subsection{Self-Interaction-Corrected Kinetic Functional}

There are arguments in the literature to approximate exact non-interacting electronic kinetic energy as a sum of the von Weizsäcker functional and the rest
\begin{equation}
    T_{\text{s}}[n]
    =
    T_{\text{W}}[n]
    +
    \text{rest}.
\end{equation}
The role of the rest can be fulfilled by any semi-local approximate functional, however, there are multiple evidence that it is better to scale the approximate functional. Particularly, it has been shown, that for the Thomas--Fermi kinetic functional, the sum of Weizsäcker and Thomas--Fermi overestimates the energy badly, so it was proposed to follow a decomposition:
\begin{equation}
    T_{\text{s}}[n]
    =
    T_{\text{W}}[n]
    +
    \gamma(N)
    T_{\text{TF}}[n],
\end{equation}
where $\gamma(N)$ is a $N$-dependent function. The specific choice of the function can be justified by both theoretical and empirical arguments, thus, Acharya \textit{et al.} proposed a form
\begin{equation}
    \gamma_{\text{ABSP}}(N)
    =
    1
    -
    \frac{C}{N^{1/3}},
\end{equation}
where, for closed-shells, $C = C_{\text{ABSP1}} = 1.412$ or $C = C_{\text{ABSP2}} = 1.332$ when fitted to 55 neutral HF atoms or 1200 atoms and ions, respectively. Gazquez and Robles analytically derived another form given by 
\begin{equation}
    \gamma_{\text{GR}}(N)
    =
    \left(
        1
        -
        \frac{2}{N}
    \right)
    \left(
        1
        -
        \frac{C_0}{N^{1/3}}
        +
        \frac{C_1}{N^{2/3}}
    \right),
\end{equation}
where $C_0 = 1.015$ and $C_1 = 0.150$.

\subsection{Self-Interaction-Corrected Exchange Functional}

In electronic structure theory, an exchange functional serves two physical roles: the interelectronic exchange and the self-interaction. Thus, the role of the interelectronic exchange is played by an approximate exchange functional, e.g., LDA or GGA approximations. 

Approximate functionals can be shown to contain a spurious approximation to the self-interaction part, which is exactly captured by FA. Therefore, we propose to describe an exchange part in LPS theory by a hybrid of the Fermi--Amaldi term and a fraction of the approximate exchange excluding its self-interaction part.

\section{Discussion}

\subsection{Spin-Dependence of Exact and Approximate Functionals}

Any approximate or exact functional can be derived for particles of arbitrary spin by considering the spin-statistics of the system. An important example is the exact exchange functional expressed in terms of one-particle reduced density matrix (1-RDM) $\rho(\mathbf{r};\mathbf{r}') = \sum_{i=1}^{N/M} \phi_i^*(\mathbf{r}) \phi_i(\mathbf{r}')$, where $\phi_i(\mathbf{r})$ is a spatial orbital, as
\begin{equation}
    E_{\text{x}}^{\text{exact}}
    [\rho; M]
    =
    -\frac{1}{2M}
    \iint
    \frac{|\rho(\mathbf{r};\mathbf{r}')|^2}
    {|\mathbf{r}-\mathbf{r}'|}
    d\mathbf{r}d\mathbf{r}'.
    \label{eq: KS Exact Exchange}
\end{equation}
Here, the explicit spin-dependence of the functional is consequential to the spin-dependence of the wave function from which the 1-RDM is derived. For fermions with $M=N$, Eq.~\eqref{eq: KS Exact Exchange} reduces to the Fermi--Amaldi (FA) exchange functional~\cite{Fermi.1934.MRAI.6.117}
\begin{equation}
    E_{\text{x}}^{\text{FA}}
    [\rho]
    =
    -\frac{1}{2N}
    \iint
    \frac{\rho(\mathbf{r})\rho(\mathbf{r}')}{|\mathbf{r}-\mathbf{r}'|}
    d\mathbf{r}d\mathbf{r}'.
    \label{eq: FA Functional}
\end{equation}
Another similar example is the non-interacting kinetic energy functional 
\begin{equation}
    T_{\text{s}}[\rho; M]
    =
    \frac{M}{2}
    \sum^{N/M}_{i=1}
    \int
    |\nabla \phi_i(\mathbf{r})|^2
    d\mathbf{r},
\end{equation}
which reduces to the von Weizsäcker (vW) kinetic energy functional for fermions with $M=N$~\cite{Weizsäcker.1935.ZP.96.431,Parr.1989.density}
\begin{equation}
    T_{\text{vW}}[\rho]
    =
    \frac{1}{8}
    \int
    \frac{|\nabla \rho(\mathbf{r})|^2}{\rho(\mathbf{r})}
    d\mathbf{r}.\label{eq: vW Functional}
\end{equation}
The Fermi--Amaldi functional is the exact exchange for one-electron systems and equivalent to Eq.~\eqref{eq: KS Exact Exchange} for two-electron singlets~\cite{Ayers.2005.MP.103.2061, Bosko.2023.JCP.159.131101}. This parallels the von Weizsäcker kinetic energy functional~\eqref{eq: vW Functional}, which is exact in these same limits but fails for many-electron systems~\cite{Weizsäcker.1935.ZP.96.431, Parr.1989.density}. 

The $s$-dependence of approximate functionals, such as LDA exchange or Thomas--Fermi kinetic energy, can be shown to arise from their statistical derivations~\cite{Parr.1989.density, Bosko.2023.PRA.108.052803}
\begin{gather}
    T_{\text{TF}}[\rho; M]
    =
    \left(\frac{2}{M}\right)^{2/3}
    C_{\text{F}}
    \int
    \rho^{5/3}(\mathbf{r})
    d\mathbf{r},\\[10pt]
    E_{\text{x}}^{\text{LDA}}[\rho; M]
    =
    \left(
    \frac{2}{M}
    \right)^{1/3}
    C_{\text{x}}
    \int
    \rho^{4/3}(\mathbf{r})
    d\mathbf{r},
\end{gather}
leading to explicit dependence on electron number for fermions with $M=N$.

\subsection{Explicit Dependence on Electron Number of Approximate Functionals}

Consider an orbital-free energy functional
\begin{equation}
    E[n]
    =
    T_{\text{W}}[n]
    +
    T_{\text{P}}[n]
    +
    V_{\text{ext}}[n]
    +
    J[n]
    +
    E_{\text{xc}}[n],
    \label{eq: OFDFT Energy}
\end{equation}
where $T_{\text{P}}[n]$ and $E_{\text{xc}}[n]$ must be approximated. For one-electron systems, $T_{\text{P}}[n]$ must vanish and $E_{\text{xc}}[n]$ must reduce to $E_{\text{x}}^{\text{FA}}$ to recover the exact energy. For a two-electron singlet the same requirements are restoring the HF energy. For higher electron numbers, these constraints no longer apply, but it is still desirable that approximate functionals reduce to the correct limits for small $N$.

One possible approximation for the exchange energy is the spin-polarized LDA exchange functional
\begin{equation}
    E_{\text{x}}^{\text{LDA}}[n_{\alpha}, n_{\beta}]
    =
    2^{1/3}
    C_{\text{x}}
    \int
    \left[
    n^{4/3}_{\alpha}(\mathbf{r}) + n^{4/3}_{\beta}(\mathbf{r})
    \right]
    d\mathbf{r},
    \label{eq: LDA Exchange}
\end{equation}
where $C_{\text{x}}=-3/4(3/\pi)^{1/3}$. This exchange functional is known to depend on the spin-statistics of the particles it is describing through its prefactor
\begin{equation}
    E_{\text{x}}^{\text{LDA}}[n; M]
    =
    \left(
    \frac{2}{M}
    \right)^{1/3}
    C_{\text{x}}
    \int
    n(\mathbf{r})^{4/3}
    d\mathbf{r},
\end{equation}
where $M = 2s+1$ and $s$ is the spin quantum number. For boson-like fermions with $M=N$, the prefactor becomes $(2/N)^{1/3}C_{\text{x}}$. Subtracting this from the original LDA exchange functional yields the $N$-dependent correction
\begin{equation}
    E_{\text{x}}^{\text{LDA}}[n_{\sigma}; N_{\sigma}]
    =
    2^{1/3}
    \left(
    1
    -
    \frac{1}{N^{1/3}_{\sigma}}
    \right)
    C_{\text{x}}
    \int
    n_{\sigma}(\mathbf{r})^{4/3}
    d\mathbf{r},
\end{equation}
where $n_{\sigma}$ is the spin density for spin component $\sigma$ and the subtracted term is treated as a self-interaction contribution. This correction vanishes in the large-$N$ limit and cancels out the LSDA exchange energy completely for one- and two-electron singlets, yet for two-electron triplets this is not true. 

\subsection{Generalized Spin-Scaling Relations for Exact and Approximate Functionals}

Unlike statistical approximations such as Thomas--Fermi or Slater exchange, both FA and vW maintain their functional forms for particles of arbitrary spin~\cite{Bosko.2023.PRA.108.052803}. Applying spin-scaling relations~\cite{Oliver.1979.PRA.20.397, Bosko.2023.PRA.108.052803}, the spin-polarized FA and vW functionals are expressed as sums over the spin-density components
\begin{gather}
    E_{\text{x}}^{\text{FA}}
    [\rho_{\alpha}, \rho_{\beta},\dots]
    =
    \sum^M_{\sigma=1}
    E_{\text{x}}^{\text{FA}}
    [\rho_{\sigma}]\\
    T_{\text{vW}}
    [\rho_{\alpha}, \rho_{\beta},\dots]
    =
    \sum^M_{\sigma=1}
    T_{\text{vW}}
    [\rho_{\sigma}],
\end{gather}
these expressions represent the spin-resolved generalizations of Eq.~\eqref{eq: FA Functional} and Eq.~\eqref{eq: OF Ts}, respectively. In contrast, statistical approximations scale with the spin degeneracy $M$ according to
\begin{gather}
    E_{\text{x}}^{\text{LDA}}
    [\rho_{\alpha}, \rho_{\beta},\dots]
    =
    M^{1/3}
    \sum^M_{\sigma=1}
    E_{\text{x}}^{\text{LDA}}
    [\rho_{\sigma}]\\
    T_{\text{TF}}
    [\rho_{\alpha}, \rho_{\beta},\dots]
    =
    M^{2/3}
    \sum^M_{\sigma=1}
    T_{\text{TF}}
    [\rho_{\sigma}].
\end{gather}
This scaling behavior applies to most approximations on the DFT Jacob's ladder, where spin dependence enters either via explicit prefactors (as in Slater exchange or Thomas--Fermi) or through variables such as the Fermi wave vector, $k_{\text{F}}$~\cite{Bosko.2023.PRA.108.052803, Parr.1989.density}.

Following Finzel \textit{et al.}~\cite{Finzel.2018.APCS.34.650}, we partition the total orbital-free energy~\eqref{eq: OFDFT Energy} into bosonic ($E_{\text{B}}$) and fermionic ($E_{\text{F}}$) components. Bosonic functionals can be derived by substituting the normalized square root of the electron density into the corresponding orbital-dependent expressions. This yields the decomposition $E[n]=E_{\text{B}}[n]+E_{\text{F}}[n]$, where $E_{\text{B}}[n]=T_{\text{vW}}[n]+V_{\text{ext}}[n]+U[n]+E_{\text{x}}^{\text{FA}}$ and $E_{\text{F}}[n]\equiv G[n]$. The bosonic part comes out of the derivation, while the fermionic part includes approximate corrections to the boson-like kinetic and particle-particle interaction parts. 

This decomposition suggests that approximations to the functional $G[n]$ are to be developed as post-boson-like energy terms, without necessarily splitting them into kinetic and exchange-correlation contributions. Consequently, while the Fermi--Amaldi functional is the exact orbital-free exchange, it constitutes only the part of the interacting electrons exchange. Improvements of this aspect are under study.

\subsection{Self-Interaction Corrected Energy Functionals}

Consider the energy functional of Eq.~\eqref{eq: OFDFT Energy} for one-electron systems. In this case, the approximate functional $G[n]$ must reduce to the FA functional to restore the exact energy. For two-electron singlets, the HF energy is exactly recovered when $G[n]=E_{\text{x}}^{\text{FA}}[n]$. Thus, the `bosonic' contribution is treated exactly by the sum of the Weizsäcker, classical Coulomb and the Fermi--Amaldi functionals since they are the expectation values of the corresponding operators. For systems with more than two electrons, both the Pauli and the KS exchange-correlation functionals are required to restore the exact energy. 

Conversely, when an approximate functional yields a non-zero value for one-electron systems or two-electron singlets, it introduces self-interaction error. This error can be eliminated by subtracting the `bosonic' contribution, which is derived from the functional's dependence on the spin statistics of the reference system~\cite{Bosko.2023.PRA.108.052803}. Specifically, for fermions with spin degeneracy $M=N$, the explicit $M$-dependence of the functional becomes a dependence on the total particle number $N$.

As an example, consider the Dirac exchange functional (LDA)~\eqref{eq: LDA exchange}. Its spin-dependent form for fermions with degeneracy $M$ is obtained by scaling the standard electronic functional of Eq.~\eqref{eq: LDA exchange} by a factor of $(2/M)^{1/3}$~\cite{Bosko.2023.PRA.108.052803}. The bosonic limit corresponds to the case where $M=N$. Consequently, the self-interaction corrected (SIC) spin-polarized~\cite{Oliver.1979.PRA.20.397,Bosko.2023.PRA.108.052803} LDA exchange functional is given by
\begin{gather}\label{eq: LDA-SIC}
    E^{\text{LDA-SIC}}_{\text{x}}[n]
    =
    2^{1/3}
    \sum_{\sigma=\alpha,\beta}
    \left(
        1 - \frac{1}{N_{\sigma}^{1/3}}
    \right)
    E^{\text{LDA}}_{\text{x}}[n_{\sigma}],
\end{gather}
where $n_{\sigma}$ is the spin density for spin component $\sigma$. This functional vanishes for one-electron systems and two-electron singlets, yet correctly yields a non-zero value for two-electron triplets.

\begin{table}[htb]
\caption{SCF atomic energies (Hartree) within the 31 $s$-type basis set. MAEs and rMAEs are given with respect to the UHF/UGBS.}
\begin{ruledtabular}
\begin{tabular}{lrrrrr}
Atom & Eq.~\eqref{eq: TFD-SIC} & Eq.~\eqref{eq: WFA} & UHF \\[2pt]
\hline\\[-9pt]
H     & $-$0.5000   & $-$0.5000   & $-$0.5000   \\[1.5pt]
He    & $-$2.8617   & $-$2.8617   & $-$2.8617   \\[1.5pt]
Li    & $-$7.0356   & $-$8.5470   & $-$7.4328   \\[1.5pt]
Be    & $-$13.2319  & $-$19.0194  & $-$14.5730  \\[1.5pt]
B     & $-$22.1259  & $-$35.7426  & $-$24.5293  \\[1.5pt]
C     & $-$33.9335  & $-$60.1803  & $-$37.6900  \\[1.5pt]
N     & $-$48.7991  & $-$93.7963  & $-$54.4045  \\[1.5pt]
O     & $-$65.9537  & $-$138.0543  & $-$74.8140  \\[1.5pt]
F     & $-$86.4703 & $-$194.4182 & $-$99.4114  \\[1.5pt]
Ne    & $-$110.6741 & $-$264.3516 & $-$128.5471 \\[1.5pt]
\hline\\[-9pt]
MAE & 5.32 & 37.27 & -- \\[1.5pt]
rMAE & 8.34 & 50.90 & -- \\
\end{tabular}
\end{ruledtabular}
\end{table}

The same logic can be applied to any approximate Pauli or exchange-correlation functional to eliminate self-interaction errors in one-electron systems and two-electron singlets. For instance, the Thomas--Fermi kinetic energy functional~\eqref{eq: TF} can be corrected as
\begin{equation}\label{eq: TF-SIC}
    T^{\text{TF-SIC}}_{\text{P}}[n]
    =
    2^{2/3}
    \sum_{\sigma=\alpha,\beta}
    \left(
        1 - \frac{1}{N_{\sigma}^{2/3}}
    \right)
    T_{\text{TF}}[n_{\sigma}].
\end{equation}
The corrections of this type vanish in large-$N$ limit restoring the original approximate functional. The resulting functionals~\eqref{eq: LDA-SIC} and \eqref{eq: TF-SIC} are resembling the functionals of~\cite{Acharya.1980.PNAS.77.6978,Cedillo.1986.JCP.85.7188} that were obtained by fitting the corresponding approximate functionals to the post-SCF atomic energies.

To implement the LDA-SIC and TF-SIC approximations for the exchange-correlation and Pauli terms, respectively, we solved Eq.~\eqref{eq: OFDFT Differential Equation} using the functional derivatives of Eqs.~\eqref{eq: LDA-SIC} and \eqref{eq: TF-SIC}, treating $N$ as a parameter~\cite{Bosko.2025.JCP.163.054108,Tozer.1997.PRA.56.2726}. Consequently, the resulting potentials are simply scaled versions of the standard LDA and TF potentials. Table III reports the SCF results for H--Ne, comparing the Eq.~\eqref{eq: OFDFT Energy} with 
\begin{equation}\label{eq: TFD-SIC}
    G[n]
    =
    E_{\text{x}}^{\text{FA}}[n]
    + 
    T^{\text{TF-SIC}}_{\text{P}}[n] 
    + 
    E^{\text{LDA-SIC}}_{\text{x}}[n]
\end{equation}
to the pure `bosonic' model with 
\begin{equation}\label{eq: WFA}
    G[n]
    =
    E_{\text{x}}^{\text{FA}}[n].
\end{equation}
All calculations employed the 31 s-type basis set and were compared against UHF/UGBS reference values.

Although the proposed TFD-SIC local functional does not surpass the accuracy of the optimal-$\lambda$ results in Table I, it exhibits superior convergence properties. Specifically, it permits convergence to tighter thresholds, e.g., 1.0E-8. 

\bibliography{../bibliography}
\end{document}

